\chapter{Source crosswalk and evidence}

\section{The consolidation map}

\begin{longtable}{L{0.12\textwidth}L{0.24\textwidth}L{0.54\textwidth}}
\toprule
Sources & Shared block & Retained mathematical work\\
\midrule
\endhead
I--IV & Arithmetic packets & Coefficient and Bernoulli refinement; cyclotomic germs and invariants; reciprocal L-value transfers; positive squarefree measure, moments, resolvent, finite Jacobi models, Gamma/Laguerre scaling; signed continuation.\\
IV--VI & Shape packets & Depth-time Gram lift; joint Gamma--sech limit; exact log-Gamma window; Fourier/Wigner coordinates; orthogonal bases; shape convolution, refinement, transfers, coherent smoothing.\\
VI--VIII & Recoupling & Positive discrete-series ladder; lowest-weight projectors; Casimir atlases; ordinary and q-deformed overlaps; measurement kernels; deformation generator; pentagon and finite compilation.\\
I, II, V & Open arithmetic boundary & One principal signed coherence target, retaining the three original conjecture-record links.\\
\bottomrule
\end{longtable}

This manuscript selects a mathematical spine rather than reproducing every invariant or example. Volume I's Wronskian negative control, Volume II's Stieltjes/resultant details, the full Wigner formulas in V, and the complete transfer/frame estimates in VI remain available in the preserved papers. The source crosswalk groups related material; it does not certify that differently stated claims are equivalent or count distinct results.

The eight originals contain 431 PDF pages. The source package also preserves the earlier special-function registry as context. It is not silently merged with the numbered series. The complete reader's machine-readable index gives each paper's physical page range, which differs from its original printed page labels.

\section{Fresh test run}

\begin{center}
\begin{tabular}{lrrr}
\toprule
Volume & Fresh pytest passed & Archived built-in passed & Source PDF pages\\
\midrule
I & 13 & 51 & 60\\
II & 18 & 61 & 50\\
III & 19 & 34 & 38\\
IV & 18 & 18 & 55\\
V & 19 & 19 & 52\\
VI & 21 & 24 & 54\\
VII & 20 & 13 & 60\\
VIII & 20 & 23 & 62\\
\midrule
Total & 148 & 243 & 431\\
\bottomrule
\end{tabular}
\end{center}

All 148 collected pytest cases passed, with zero failures, errors, or skips. Each suite ran from a clean temporary source copy. SHA256 comparisons confirmed that the original extracted source trees remained unchanged. Tests include recomputed numerical and symbolic identities as well as archival-ledger, figure, and manuscript checks.

The archived built-in ledgers report 243 successful checks. Only the 64 built-in checks in Papers I and VII were recomputed by the current pytest run. The remaining built-in totals are historical records. Neither test count is a universal proof count. The JSON report preserves exact commands, environment versions, timings, and case-level results. Small exact correction witnesses are recorded separately.

\section{Formalization scope}

The retrieved projector-calculus ledger has 173 special-function source objects: 125 theorem, proposition, lemma, or corollary records; 15 definitions; 22 context, example, or algorithm records; eight boundary records; and three conjecture records.

\begin{center}
\begin{tabular}{lrr}
\toprule
Volume & Source objects & Theorem/lemma/proposition/corollary records\\
\midrule
I & 32 & 22\\
II & 23 & 15\\
III & 24 & 16\\
IV & 16 & 10\\
V & 21 & 17\\
VI & 22 & 20\\
VII & 15 & 10\\
VIII & 20 & 15\\
\midrule
Total & 173 & 125\\
\bottomrule
\end{tabular}
\end{center}

The statuses are 88 phase-1-queued, 51 queued, one Lean-draft-present, 22 context/example, eight boundary, and three conjectural. No special-function claim in that retrieved ledger has a kernel-checked source-claim status. The actual special-function Lean tranche inspected contains one coefficient-subsequence refinement theorem. It does not formalize the complete analytic series, Gamma limits, orthogonal-polynomial analysis, or signed coherence estimate.

The whole-project v1.1.0 static figures of 594 claim records and 104 theorem-like Lean declarations are not special-function proof counts. That release explicitly reports that a kernel build was not attempted. The later finite browser-safe core does not establish analytic completion of this series. These findings are bounded by the targeted files and searches recorded in the scope report.

\section{Primary-source checks}

The following exact external inputs received bounded primary-source checks in this edition. Other source-paper bibliographies remain preserved for subsequent normalization.

\begin{enumerate}[leftmargin=*]
\item L. B\'aez-Duarte, \emph{A sequential Riesz-like criterion for the Riemann hypothesis}, International Journal of Mathematics and Mathematical Sciences 2005, 3527--3537, Theorem 1.1. DOI: \href{https://doi.org/10.1155/IJMMS.2005.3527}{10.1155/IJMMS.2005.3527}. Its binomial coefficients agree with Chapter 1's $c_k$ by the exact reciprocal-zeta identity; its arbitrary-positive-exponent bound gives the stated coherence equivalence.
\item NIST DLMF, \href{https://dlmf.nist.gov/18.28}{Section 18.28}, equations 18.28.1\_5 and 18.28.6\_1--6\_2. These match the normalized scalar Askey--Wilson function and the q-difference operator in VIII. The separate orthogonality parameter restrictions are not inferred from the eigenfunction identity.
\item W. Gautschi, \emph{Orthogonal polynomials: applications and computation}, Acta Numerica 5 (1996), 45--119, opening Gaussian-quadrature discussion and support-interval statements. The quadrature support is an interval containing the measure's support, not an assertion that quadrature nodes are original atoms. The decisive counterexample in Chapter 5 is proved directly from the packet measure.
\end{enumerate}

The accompanying literature notes state exactly what was checked and what was not. The consolidation makes no global novelty claim. Classical special functions retain their established names; the unified organization records how reconstruction, refinement, positive models, shape transfer, and recoupling are composed within the packet calculus.
